\documentclass[10pt,a4paper]{amsart}
\usepackage[margin=19mm]{geometry}
\usepackage{amsmath,amssymb,mathtools}
\usepackage[scr=rsfs,cal=euler]{mathalfa}
\usepackage{xfrac}
\usepackage[unicode,hidelinks,bookmarks=false]{hyperref}
\newcommand{\ie}{i.e.\ }
\newcommand{\cf}{cf.\ }
\newcommand{\resp}{respectively\ }
\newcommand{\Subsection}{\subsection}
\providecommand{\nobreakdash}{\nobreak\hbox{-}}
\setlength{\emergencystretch}{2em}
\multlinegap0pt
\usepackage{amssymb}
\usepackage{xcolor}
\newtheorem{theorem}{Theorem}
\title{The radial Newton problem: nonlinear dynamics of minimal resistance in central fields}
\author{Rafael L\'opez}
\date{Source: arXiv:2605.14383v1 (CC BY 4.0)}
\begin{document}
\maketitle

\noindent\emph{Source and licence.} The radial Newton problem: nonlinear dynamics of minimal resistance in central fields. Version arXiv:2605.14383v1 (CC BY 4.0), distributed by arXiv under the Creative Commons Attribution 4.0 International licence, http://creativecommons.org/licenses/by/4.0/, as recorded in the arXiv licence metadata for this exact version and shown on the abstract page https://arxiv.org/abs/2605.14383v1. Full licence notice: \texttt{licenses/arxiv-2605.14383-CC-BY-4.0.txt}.\par\smallskip

\noindent\textit{Editorial scope.} Counted unit: the article's theorem t38 (source0.tex lines 704--707). The article's complete original proof of it is delivered with it (source0.tex lines 709--776, one \texttt{proof} environment of 68 lines). No mathematics is written, repaired or re-derived: the statement and the proof are the article's own lines, in the article's own notation, and no retained line is substituted or re-flowed.  Retained uncounted as the source-local context the proof needs: lines 697--699.  Nothing was omitted from any retained span, so the package carries no omission record.  No retained line contains a citation, so the wrapper carries no bibliography and no bibliography entry.  No retained line contains a figure environment, an image-inclusion command or drawing code, so no image is delivered and the package declares no asset dependency.  Editorial formatting only, in addition: an AMS \texttt{amsart} preamble that restates the article's own declarations and macros used by the retained lines, namely the environments \texttt{theorem} on separate counters, together with the standard packages those lines need.  The retained environments print in this document as Theorem 1 in order of appearance, whereas the article prints the same items with the numbering of its own full document.  The complete native source of the article as distributed in the arXiv e-print for this exact version is bundled unchanged as \texttt{sources/200579/source0.tex}.
\begin{equation}\label{m23}
 \frac{d}{d\theta} \left( \frac{\sin\theta \, u'}{(1 + u'^2)^2} \right) = \sin\theta \frac{u'^2 - 1}{(1 + u'^2)^2}.
\end{equation}

\begin{theorem} \label{t38}
For any initial height $u_0 \in \mathbb{R}$, there exists a $\delta > 0$ and a unique strictly concave, smooth function $u \in C^2([0, \delta))$ that satisfies the Euler-Lagrange equation \eqref{m23} with $u(0) = u_0$ and $u'(0) = 0$. Furthermore, $u''(0) = -\frac12$ and the solution $u$ has the following local expansion:
$$u(\theta) = u_0 - \frac{1}{4}\theta^2 + O(\theta^4).$$
\end{theorem}

\begin{proof}
Expanding the Euler-Lagrange equation yields
\begin{equation} \label{sin}
 u'' + \cot\theta \, u' \left( \frac{1+u'^2}{1-3u'^2} \right) = \frac{(1+u'^2)(u'^2-1)}{1-3u'^2} .
\end{equation}
This is a second-order nonlinear ODE with a singularity at $\theta = 0$. We require the solution to be smooth at $\theta=0$, hence $u'(0) = 0$. Assuming the solution $u$ reaches $\theta=0$, by taking the limit as $\theta \to 0$ in \eqref{sin}, the right-hand side tends to $-1$. On the left-hand side, as $\theta \to 0$, the first term tends to $u''(0)$. 
For the second term, by L'H\^{o}pital's rule, $\lim_{\theta \to 0} u'(\theta)\frac{\cos\theta}{\sin\theta} = u''(0)$. Therefore, 
the entire left-hand side evaluates to $2u''(0)$. Equating this to 
the right-hand side limit, which is $-1$, we conclude $u''(0) =- \frac12$.
 
 To establish the existence and uniqueness of such a solution, we rewrite \eqref{sin} in the form of a first-order system for the variables $(u, v)$ where $v = u'$. The equation can be expressed as:
\begin{equation} \label{prime}
 v' + \frac{v}{\theta} F(\theta, v) = G(v),
\end{equation}
where 
$$F(\theta, v) = \theta \cot \theta \left( \frac{1+v^2}{1-3v^2} \right),\quad G(v) = \frac{(1+v^2)(v^2-1)}{1-3v^2}.$$
 Note that $F$ and $G$ are analytic in a neighborhood of $\theta = 0, v = 0$, with $F(0, 0) = 1$ and $G(0) = -1$. The singularity of $\cot \theta$ at the origin is removable in the product $\theta \cot \theta$, making $F$ analytic in its domain.


Multiplying \eqref{prime} by $\theta$ and rearranging, we can write the equation as:
\begin{equation*}
 (\theta v)' = \theta \left[ G(v) + \frac{v}{\theta} (1 - F(\theta, v)) \right].
\end{equation*}
Integrating from $0$ to $\theta$ and dividing by $\theta$, we obtain the fixed-point integral equation for $v(\theta)$:
\begin{equation} \label{ope}
 v(\theta) = \frac{1}{\theta} \int_0^\theta s \left[ G(v(s)) + \frac{v(s)}{s} (1 - F(s, v(s))) \right] ds.
\end{equation}
 
 
Since we seek a $C^2$ solution with $u'(0)=0$ and $u''(0)=-\frac12$, we define a new change of variable $v(\theta) = \theta w(\theta)$, where $w(0)=-\frac12$. Substituting this into \eqref{ope}, we define the operator $\mathsf{T}$ on the space of continuous functions $C^0([0, \delta])$ as:
\begin{equation*}
 (\mathsf{T}w)(\theta) = \frac{1}{\theta^2} \int_0^\theta s H(s, w(s)) \, ds,
\end{equation*}
where 
$$H(\theta, w) = G(\theta w) + w(1 - F(\theta, \theta w)).$$
Since $1 - F(\theta, \theta w) = O(\theta^2)$ as $\theta \to 0$, the operator $\mathsf{T}$ is well-defined and continuous at the origin.  For a given $\epsilon > 0$, let $\overline{B(-\frac12, \epsilon)} \subset C^0([0, \delta])$ be the closed ball centered at the constant function $w =- \frac12$. Since $F$ and $G$ are analytic, $H$ is smooth near $(0, -\frac12)$ with $H(0, -\frac12) = -1$. Consequently, there exist constants $K, M > 0$ such that for all $\theta \in [0, \delta]$ and $w, w_1, w_2 \in \overline{B(-\frac12, \epsilon)}$:
$$
|H(\theta, -\frac12) +1| \leq M\theta, \quad |H(\theta, w_1) - H(\theta, w_2)| \leq K \theta |w_1 - w_2|.
$$
By the triangle inequality, for any $w \in \overline{B(-\frac12, \epsilon)}$, we obtain the bound 
$$|H(\theta, w) +1| \leq |H(\theta, w) - H(\theta, -\frac12)| + |H(\theta, -\frac12) +1| \leq (K\epsilon + M)\theta.$$
 We perform the proof in two steps.

\begin{enumerate}
 \item Invariance. We prove that $\mathsf{T}(\overline{B(-\frac12, \epsilon)}) \subset \overline{B(-\frac12, \epsilon)}$. Indeed, using the bound derived above,
 \begin{equation*}
 \begin{split}
 |(\mathsf{T}w)(\theta) +\frac12| &\leq \frac{1}{\theta^2} \int_0^\theta s |H(s, w(s)) +1| \, ds \leq \frac{1}{\theta^2} \int_0^\theta s^2 (K\epsilon + M) \, ds \\
 &= (K\epsilon + M)\frac{\theta}{3} \leq (K\epsilon + M)\frac{\delta}{3}.
 \end{split}
 \end{equation*}
 Choosing $\delta$ small enough such that  $(K\epsilon + M)\frac{\delta}{3} \leq \epsilon$ ensures the ball is invariant under the operator $\mathsf{T}$.
 
 \item Contraction.   Given $w_1, w_2 \in \overline{B(-\frac12, \epsilon)}$, we have:
 \begin{equation*}
 \begin{split}
 |(\mathsf{T}w_1)(\theta) - (\mathsf{T}w_2)(\theta)| &\leq \frac{1}{\theta^2} \int_0^\theta s |H(s, w_1(s)) - H(s, w_2(s))| \, ds \\
 &\leq \frac{K \|w_1 - w_2\|_\infty}{\theta^2} \int_0^\theta s^2 \, ds = \frac{K \theta}{3} \|w_1 - w_2\|_\infty.
 \end{split}
 \end{equation*}
 For $\frac{K R}{3} < 1$, the operator $\mathsf{T}$ is a contraction.
\end{enumerate}
 
 
 By the Banach fixed-point theorem, there exists a unique $w \in C^0([0, R])$. Then $u(\theta) = u_0 + \int_0^\theta s w(s) \, ds$ is the unique $C^2$ solution satisfying the initial conditions.  Finally, since $u \in C^2$ and $u''(0) = -\frac12 < 0$, we can choose $\delta \leq R$ sufficiently small such that $u''(\theta) < 0$ on $[0, \delta)$, guaranteeing that $u(\theta)$ is strictly concave. The Taylor expansion follows directly from the initial conditions.

 
\end{proof}

\end{document}
